\documentclass[10pt,a4paper]{amsart}
\usepackage[margin=19mm]{geometry}
\usepackage{amsmath,amssymb,mathtools}
\usepackage[scr=rsfs,cal=euler]{mathalfa}
\usepackage{xfrac}
\usepackage[unicode,hidelinks,bookmarks=false]{hyperref}
\newcommand{\ie}{i.e.\ }
\newcommand{\cf}{cf.\ }
\newcommand{\resp}{respectively\ }
\newcommand{\Subsection}{\subsection}
\providecommand{\nobreakdash}{\nobreak\hbox{-}}
\setlength{\emergencystretch}{2em}
\multlinegap0pt
\usepackage{amssymb}
\usepackage{dsfont}
\newtheorem{proposition}{Proposition}
\newcommand{\Cnorm}[1]{\left\|#1\right\|}
\DeclareMathOperator{\E}{\mathbb{E}}
\newcommand{\R}{\ensuremath{\mathbb{R}}}
\newcommand{\abs}[1]{\left|#1\right|}
\newcommand{\chibar}{\overline \chi}
\newcommand{\frcmod}[1]{\left\Vert#1\right\Vert}
\renewcommand{\mod}[1]{{\ifmmode\text{\rm\ (mod~$#1$)}\else\discretionary{}{}{\hbox{ }}\rm(mod~$#1$)\fi}}
\DeclareMathOperator{\prob}{Prob}
\title{Distribution of mixed character sums and extremal problems for Littlewood polynomials}
\author{Jonathan W. Bober, Oleksiy Klurman, Besfort Shala}
\date{Source: arXiv:2510.06161v2 (CC BY 4.0)}
\begin{document}
\maketitle

\noindent\emph{Source and licence.} Distribution of mixed character sums and extremal problems for Littlewood polynomials. Version arXiv:2510.06161v2 (CC BY 4.0), distributed by arXiv under the Creative Commons Attribution 4.0 International licence, http://creativecommons.org/licenses/by/4.0/, as recorded in the arXiv licence metadata for this exact version and shown on the abstract page https://arxiv.org/abs/2510.06161v2. Full licence notice: \texttt{licenses/arxiv-2510.06161-CC-BY-4.0.txt}.\par\smallskip

\noindent\textit{Editorial scope.} Counted unit: the article's proposition proposition:truncation (source0.tex lines 832--849). The article's complete original proof of it is delivered with it (source0.tex lines 850--1017, one \texttt{proof} environment of 168 lines). No mathematics is written, repaired or re-derived: the statement and the proof are the article's own lines, in the article's own notation, and no retained line is substituted or re-flowed.  Retained uncounted as the source-local context the proof needs: lines 694--721.  Nothing was omitted from any retained span, so the package carries no omission record.  No retained line contains a citation, so the wrapper carries no bibliography and no bibliography entry.  No retained line contains a figure environment, an image-inclusion command or drawing code, so no image is delivered and the package declares no asset dependency.  Editorial formatting only, in addition: an AMS \texttt{amsart} preamble that restates the article's own declarations and macros used by the retained lines, namely the environments \texttt{proposition} on separate counters, together with the standard packages those lines need.  The retained environments print in this document as Proposition 1 in order of appearance, whereas the article prints the same items with the numbering of its own full document.  The complete native source of the article as distributed in the arXiv e-print for this exact version is bundled unchanged as \texttt{sources/186587/source0.tex}.
\begin{proposition}\label{prop:poisson}
    For a primitive Dirichlet character $\chi$ mod $q$ and for $\alpha,\beta \in \R$,
    we have
\begin{multline*}
    F_{k,\chi,\alpha,\beta}(t)
    = e(\alpha t)\frac{\tau(\chi)}{2\pi i q^{1/2}}\sum_{\abs{l} < K}
    \frac{e(\alpha l) (e(\beta(l + t)) - 1)}
    {l + t }
    \chibar(k-l)
    \\
    +
    O\left(\frac{q^{1/2} \log(q)}{K} + \frac{1}{q^{1/2}}\min\left(1, \frac{1}{K \frcmod{\alpha q}}\right)
        + \frac{1}{q^{1/2}}\min\left(1, \frac{1}{K\frcmod{(\alpha + \beta)q}}\right)
    + \frac{\log K}{Kq^{1/2}}
    \right).
\end{multline*}
    In particular, for any $\alpha$ and $\beta$ such that $\alpha q$ and
    $(\alpha + \beta)q$ are not integers, we have
    \begin{equation}\label{elegant_formula}
    F_{k,\chi,\alpha,\beta}(t) =
    e(\alpha t)\frac{\tau(\chi)}{2\pi i q^{1/2}}
    \lim_{K \rightarrow \infty}
    \sum_{\abs{l} < K}
    \frac{e(\alpha l) (e(\beta(l + t)) - 1)}
    {l + t }
    \chibar(k-l).
\end{equation}
\end{proposition}

\begin{proposition}\label{proposition:truncation}
    Let $L(q)$ be some function such that $L(q) \rightarrow \infty$ and
    $L(q) \ll q^\varepsilon$ for all $\varepsilon > 0$.
    For fixed real numbers $\alpha$ and $\beta$, any bounded Lipschitz function $h:\mathscr{C}[0,1] \rightarrow \R$, and
    any sequence of integers $k_q$, as $q$ tends to infinity over the primes we have 
    %the expected value $\E h(F_{k_q,q,\alpha,\beta})$
    %approaches $\E h(\tilde F_{k_q,q,\alpha,\beta})$ as $q$ tends to infinity over the
    %primes, in the sense that
    \[
        \abs{\E h(F_{k_q,q,\alpha,\beta}) - \E h(\tilde F_{k_q,q,\alpha,\beta})} \rightarrow 0 \text{ and } \abs{\E h(G_{q,\alpha,\beta}) - \E h(\tilde G_{q,\alpha,\beta})} \rightarrow 0.
    \]
    That is, $F_{k_q,q,\alpha,\beta}$ converges weakly to some $F$ (as $q$ tends
    to infinity over the primes, or over some subsequence of the primes) if and
    only if $\tilde F_{k_q,q,\alpha,\beta}$ converges weakly to $F$, and similarly for $G_{q, \alpha, \beta}$ and $\tilde G_{q, \alpha, \beta}$. 

    %Similarly, the expected value $\E h(G_{q,\alpha,\beta})$ approaches
   % $\E h(\tilde G_{q,\alpha,\beta})$.
\end{proposition}

\begin{proof}
    We will abbreviate $F_{k,\chi} := F_{k,\chi,\alpha,\beta}$.
    By the triangle inequality
    \begin{multline*}
        \abs{F_{k,\chi}(t) - \tilde F_{k,\chi}(t)}^2
        \ll \abs{\tilde F_{k,\chi}(t) -
            \frac{\tau(\chi)}{2 \pi i q^{1/2}} \sum_{\abs{n} < q}
    \frac{e(\alpha n)(e(\beta(n + t) - 1)}
        {t + n}\chibar(k - n)
            }^2 \\
            +
        \abs{F_{k,\chi}(t) -
            \frac{\tau(\chi)}{2 \pi i q^{1/2}} \sum_{\abs{n} < q}
    \frac{e(\alpha n)(e(\beta(n + t) - 1)}
        {t + n}\chibar(k - n)
            }^2.
    \end{multline*}
    The second term is $O((\log q)/q^{1/2})$ by the quantitative bound of
    Proposition \ref{prop:poisson} for any nontrivial character $\chi$ mod $q$,
    so
    \[
    \abs{F_{k,\chi}(t) - \tilde F_{k,\chi}(t)}^2
        \ll
            \abs{
            \sum_{L(q) < |n| < q}
            \frac{e(\alpha n)(e(\beta(n + t) - 1)}
            {t + n}
            \chibar(k - n)
        }^2 + \frac{\log q}{q}.
    \]
    For $t \in [0,1]$ and $|n| > L(q)$, we have $1/(n-t) = 1/n + O(1/n^2)$,
    so again from the triangle inequality we have
    \begin{multline}\label{eq:supbound}
        \sup_{t\in[0,1]} \abs{F_{k,\chi}(t) - \tilde F_{k,\chi}(t)}^2 \ll
        \\
            \abs{
            \sum_{L(q) < |n| < q}
            \frac{e((\alpha + \beta)n)}
            {n}
            \chibar(k - n)
        }^2 +
            \abs{
            \sum_{L(q) < |n| < q}
            \frac{e(\alpha n)}
            {n}
            \chibar(k - n)
        }^2
             + \frac{1}{L(q)^{2}}.
     \end{multline}
    Now taking the average over all nontrivial characters we find that
    \begin{multline*}
        \frac{1}{q-2}\sum_{\chi \ne \chi_0} \sup_{t\in[0,1]} \abs{F_{k,\chi}(t) - \tilde F_{k,\chi}(t)}^2 \ll \\
        \frac{1}{q-2}
        \sum_{\chi \mod q}
            \abs{
            \sum_{L(q) < |n| < q}
            \frac{e((\alpha + \beta)n)}
            {n}
            \chibar(k - n)
        }^2 + \\
        \frac{1}{q-2}
        \sum_{\chi \mod q}
            \abs{
            \sum_{L(q) < |n| < q}
            \frac{e(\alpha n)}
            {n}
            \chibar(k - n)
        }^2
        + \frac{1}{L(q)^{2}}
        \\
        \ll 
        \frac{q-1}{q-2}\sum_{|n| > L(q)} \frac{1}{n^2} + \frac{1}{L(q)^2}
        \ll
        \frac{1}{L(q)}.
    \end{multline*}
    Upon applying Markov's inequality it follows that the number of primitive $\chi$ such that
    $\|F_{k, \chi,\alpha,\beta} - \tilde F_{k, \chi,\alpha,\beta}\| > z$ is bounded by
    $O\left(\frac{q}{z^2 L(q)}\right)$.

    We now suppose that $h:\mathscr C[0,1] \rightarrow \R$ is a bounded Lipschitz function with \linebreak
    $\abs{h(x_1) - h(x_2)} \le K\|x_1 - x_2\|$ and $|h(x)| \le B$ for all $x_1$ and
    $x_2$. Then for any $z$ we have
    \[
    \begin{split}
        \E h(F_{k,q,\alpha,\beta}) - \E h(\tilde F_{k,q,\alpha,\beta}) &= 
        \E \left[h(F_{k, q,\alpha,\beta}) - h(\tilde F_{k, q,\alpha,\beta})\right] \\
        &\le
        Kz + B\prob\left[\Cnorm{F_{k, \chi,\alpha,\beta} - \tilde F_{k, \chi,\alpha,\beta}} > z\right] \\
        &\le Kz + O\left(\frac{B}{z^2 L(q)} + \frac{B}{q}\right)
    \end{split}
    \]
    (the term $B/q$ arises from the contribution of the trivial character).
    Choosing $z = L(q)^{-1/3}$ yields desired conclusion.

    For the quadratic case, we start with \eqref{eq:supbound} and average over
    $k$.
    Notice that
    \begin{multline*}
        \frac{1}{q}
        \sum_{k=0}^{q-1}
            \abs{
            \sum_{L(q) < |n| < q}
            \frac{e(xn)}
            {n}
            \chibar(k - n)
        }^2
        \\
        \ll
        \frac{1}{q}
        \sum_{k=0}^{q-1}
            \abs{
            \sum_{L(q) < n < q}
            \frac{e(xn)}
            {n}
            \chibar(k - n)
        }^2
        +
        \frac{1}{q}
        \sum_{k=0}^{q-1}
            \abs{
            \sum_{L(q) < -n < q}
            \frac{e(xn)}
            {n}
            \chibar(k - n)
        }^2
    \end{multline*}
    and
    \[
        \frac{1}{q}
        \sum_{k=0}^{q-1}
            \abs{
            \sum_{L(q) < \epsilon n < q}
            \frac{e(xn)}
            {n}
            \chibar(k - n)
        }^2
        =
        \frac{1}{q} \hspace{.3em}
        \sum\sum_{\hspace{-2.2em}\substack{L(q) < \epsilon m < q \\ L(q) < \epsilon n < q } }
            \frac{e(xn - xm)}
            {nm}
            \sum_{k=0}^{q-1}
            \chibar(k - n)
            \chi(k - m)
    \]
    for each choice of $\epsilon = \pm 1$.
    When $\chi$ is the quadratic character, the inner sum is $-1$ unless $n = m$,
    in which case it is $q-1$. Thus
    \[
    \begin{split}
        \frac{1}{q}
        \sum_{k=0}^{q-1}
            \abs{
            \sum_{L(q) < \epsilon n < q}
            \frac{e(xn)}
            {n}
            \chibar(k - n)
        }^2
        &=
        \frac{q-1}{q} \sum_{L(q) < \epsilon n < q} \frac{1}{n^2} -
            \frac{1}{q} \sum_{\substack{L(q) < \epsilon n < q \\ L(q) < \epsilon m < q \\ n \ne m}}
            \frac{e(xn - xm)}{nm} \\
        &\ll \frac{1}{L(q)} + \frac{(\log q)^2}{q} \ll \frac{1}{L(q)}.
        \end{split}
    \]
    We therefore have exactly the same quality bounds when we average over $k$ as when we average
    over $\chi$ and the proof is now complete.
\end{proof}

\end{document}
